\chapter{Introdução}

Sistemas embarcados estão presentes em uma ampla gama de aplicações, que vão desde dispositivos médicos e sistemas de automação industrial
até plataformas automotivas e missões espaciais. Diferentemente da computação de propósito geral, esses sistemas operam sob rigorosos requisitos
de tempo real, nos quais a corretude do comportamento não depende apenas da precisão lógica dos resultados produzidos, mas também do momento
exato em que esses resultados são disponibilizados. Em muitos casos, atrasos mínimos ou falhas temporais podem comprometer a segurança, a
confiabilidade e até mesmo a integridade física de pessoas e equipamentos.

Tradicionalmente, a garantia desse determinismo esteve associada ao uso de plataformas dedicadas, nas quais \glspl{rtos} são empregados para assegurar
previsibilidade temporal, isolamento entre tarefas e cumprimento de restrições rígidas de escalonamento. Entretanto, a evolução do software contemporâneo
deslocou grande parte dos ecossistemas de desenvolvimento — bibliotecas, frameworks, ferramentas e serviços — para ambientes baseados em \glspl{gpos} - 
especialmente Linux. Embora esses sistemas ofereçam grande riqueza funcional e amplo suporte da comunidade, sua adoção em contextos críticos embarcados demanda
esforços adicionais de adaptação, validação e integração.

Como resultado, cresce a necessidade de consolidar, em uma mesma plataforma computacional, cargas de trabalho com requisitos distintos e frequentemente
conflitantes: de um lado, aplicações críticas, sensíveis a prazos temporais e dependentes de comportamento determinístico; de outro, aplicações complexas,
orientadas a funcionalidades avançadas e dependentes de ambientes generalistas.

Nesse cenário, técnicas de virtualização surgem como uma alternativa promissora, pois permitem a coexistência de múltiplos ambientes de execução sobre o
mesmo hardware. Dessa forma, torna-se possível executar aplicações críticas e não críticas de maneira controlada, reduzindo custos de integração e
aumentando o aproveitamento da plataforma.

Paralelamente, a arquitetura aberta RISC-V vem ganhando relevância  devido à sua flexibilidade, modularidade e independência
tecnológica. Seu conjunto de instruções extensível, aliado à ausência de barreiras proprietárias, favorece tanto a pesquisa acadêmica quanto o desenvolvimento
industrial, especialmente em soluções customizadas e de baixo custo.

Por isso, este trabalho propõe o projeto e a implementação, utilizando a linguagem C++ orientada a objetos de baixo overhead, de uma infraestrutura de virtualização para sistemas
baseados em RISC-V. O foco principal reside na previsibilidade temporal, no isolamento rígido entre cargas de trabalho e no baixo custo de execução. Busca-se, assim, contribuir para
a viabilização de plataformas embarcadas capazes de integrar aplicações heterogêneas sem comprometer requisitos críticos de desempenho e confiabilidade.

\section{Objetivos}

\subsection{Objetivos Gerais}

Projetar, implementar e validar uma infraestrutura de virtualização de baixo overhead para sistemas embarcados baseados na arquitetura RISC-V, utilizando a linguagem C++,
com foco no isolamento de cargas de trabalho, no particionamento estático de recursos e na previsibilidade temporal necessária à execução de aplicações de tempo real.

\subsection{Objetivos Específicos}

\begin{enumerate}
\item Definir a arquitetura de software do hipervisor, estabelecendo os mecanismos necessários para o gerenciamento e isolamento das máquinas virtuais.
\item Implementar o núcleo de virtualização responsável pela criação, inicialização, gerenciamento e execução das máquinas virtuais em plataformas RISC-V.
\item Desenvolver mecanismos de virtualização de interrupções e exceções, assegurando o tratamento correto e eficiente dos eventos entre hardware, hipervisor e sistemas convidados.
\item Implementar mecanismos de particionamento estático e isolamento de recursos computacionais, incluindo processadores, memória e dispositivos de entrada e saída.
\item Garantir a execução isolada de aplicações com diferentes requisitos operacionais, possibilitando a coexistência de cargas de trabalho críticas e não críticas em uma mesma plataforma computacional.
\item Validar funcionalmente a infraestrutura desenvolvida por meio da execução de sistemas convidados representativos.
\item Avaliar o impacto da virtualização sobre métricas relacionadas ao desempenho e à previsibilidade temporal, incluindo latência, tempo de resposta e sobrecarga de execução.
\end{enumerate}

